\chapter{Le sette specie di settima}
\label{chap:sette_specie}

Dopo aver studiato la settima di dominante e i suoi derivati diretti, è necessario esplorare l'intera famiglia degli accordi di quadriade. In questo capitolo sistematizzeremo le diverse "specie" di settima, ovvero le diverse combinazioni di terze sovrapposte che possono nascere sui gradi della scala.

Per quanto riguarda il comportamento generale delle dissonanze, la loro preparazione e la loro risoluzione obbligatoria, si rimanda integralmente al \textbf{Capitolo \ref{chap:teoria_dissonanza}}.

\section{Classificazione generale}
\label{sec:classificazione_specie}

Tradizionalmente si distinguono sette specie di accordi di settima, determinate dalla natura della triade di base e dalla qualità dell'intervallo di settima rispetto alla fondamentale.

\begin{enumerate}
    \item \textbf{Prima specie:} Triade Maggiore + settima minore (Settima di dominante).
    \item \textbf{Seconda specie:} Triade minore + settima minore.
    \item \textbf{Terza specie:} Triade diminuita + settima minore (Settima semidiminuita).
    \item \textbf{Quarta specie:} Triade Maggiore + settima Maggiore.
    \item \textbf{Quinta specie:} Triade diminuita + settima diminuita (Settima diminuita).
    \item \textbf{Sesta specie:} Triade minore + settima Maggiore.
    \item \textbf{Settima specie:} Triade aumentata + settima Maggiore.
\end{enumerate}

In un corso di armonia complementare, ci concentreremo solo sulle casistiche più comuni e rilevanti per l'analisi e la composizione standard, ignorando le specie più rare che non trovano applicazione pratica negli esami.

\section{Prima e Terza specie: richiami}
\label{sec:richiami_specie}

Le specie più importanti sono già state trattate nei capitoli precedenti:

\begin{itemize}
    \item \textbf{Prima Specie (Settima di dominante):} Si rimanda al \textbf{Capitolo \ref{chap:settima_dominante}}. È l'unica specie che non richiede preparazione obbligatoria della settima.
    \item \textbf{Terza Specie (Settima di sensibile / semidiminuita):} Si rimanda al \textbf{Capitolo \ref{chap:vii7_ii7}}. Ricordiamo che questa specie richiede sempre la \textbf{preparazione} della settima (fase di consonanza nell'accordo precedente) e risolve scendendo di grado. La sua funzione è di dominante (sul VII grado del maggiore) o di sottodominante (sul II grado del minore).
\end{itemize}

\section{Seconda Specie}
\label{sec:seconda_specie_dettaglio}

La seconda specie è formata da una \textbf{triade minore} e una \textbf{settima minore}. L'accordo di gran lunga più importante di questa specie è quello costruito sul \textbf{secondo grado (II)} del modo maggiore.

\subsection{Funzione e risoluzione}
La sua funzione è tipicamente di \textbf{sottodominante}. Essendo una dissonanza di seconda specie, la settima deve essere rigorosamente \textbf{preparata}. Risolve naturalmente sulla dominante (V o V7).

\begin{regola}[Collegamento I-II7-V e raddoppio della terza]
    \label{reg:ii7_raddoppio}
    Nel collegamento tra il I grado e il II7 (e la successiva risoluzione su V), è frequente e consigliabile \textbf{raddoppiare la terza} dell'accordo di II7 (che corrisponde al IV grado della scala). Questo espediente tecnico è fondamentale per evitare il rischio di \textbf{quinte parallele} tra la fondamentale e la quinta del II grado nel movimento verso l'accordo di dominante.
\end{regola}

\begin{esempio}[II7 in Do Maggiore]
    \musica{
        \partilate
        { \key do \major \time 4/4 }
        { do''2 do''2 | si'1 }
        { sol'2 fa'2 | re'1 }
        { mi'2 fa'2 | fa'1 }
        { do2 re2 | sol,1 }
        \lyricmode { "I"2 "II7"2 | "V"1 }
        \figuremode { s2 <7>2 | s1 }
    }
    Si noti il raddoppio del Fa (terza del II7) tra contralto e tenore.
\end{esempio}

\section{Quarta Specie}
\label{sec:quarta_specie_dettaglio}

La quarta specie è formata da una \textbf{triade Maggiore} e una \textbf{settima Maggiore}. Nel corso di armonia studieremo solo l'accordo costruito sul \textbf{primo grado (I)} del modo maggiore.

\subsection{Uso e risoluzione}
Questo accordo ha una tensione molto particolare e risolve generalmente sulla \textbf{sottodominante (IV)}. La settima (che è la sensibile della scala) deve essere preparata e risolve scendendo di grado verso la terza del IV grado.

\begin{esempio}[I7 della quarta specie]
    \musica{
        \partilate
        { \key do \major \time 4/4 }
        { sol'2 sol'2 | la'1 }
        { mi'2 mi'2 | fa'1 }
        { do'2 si2 | do'1 }
        { do2 do2 | fa,1 }
        \lyricmode { "I"2 "I7 (4a sp.)"2 | "IV"1 }
        \figuremode { s2 <7>2 | s1 }
    }
\end{esempio}

\section{Quinta Specie}
\label{sec:quinta_specie_dettaglio}

La quinta specie è la \textbf{settima diminuita} (triade diminuita + settima diminuita). È l'accordo più dissonante e simmetrico.

\subsection{Uso sul VII grado del modo minore}
Si studia principalmente sul \textbf{VII grado del modo minore} (sulla sensibile), dove svolge funzione di dominante estrema. Come la settima di dominante, può talvolta omettere la preparazione, ma la risoluzione è obbligatoria e tutte le voci tendono a muoversi verso la tonica.

Oltre che sul VII grado, questa specie trova un uso importante come \textbf{accordo alterato sul IV grado} di entrambi i modi (maggiore e minore), agendo come una sorta di "dominante della dominante" o per scopi coloristici nelle cadenze.

\section{Sesta e Settima Specie}
\label{sec:sesta_settima_specie}

Queste due specie sono estremamente rare e non trovano uso pratico in un corso di armonia complementare. È sufficiente sapere dove si collocano teoricamente:

\begin{itemize}
    \item \textbf{Sesta specie (minore-Maggiore):} si trova sul I grado della scala minore armonica e melodica; risolve generalmente sul IV grado.
    \item \textbf{Settima specie (aumentata-Maggiore):} si trova sul III grado della scala minore armonica e melodica; risolve sul VI grado.
\end{itemize}

Nonostante la loro esistenza teorica, non verranno approfondite ulteriormente in quanto non richieste nelle prove d'esame.
